\section{Multiplication, duplication, and special values}\label{sec:multiplication}

\subsection{The general multiplication formula}

For an integer \(n\ge1\), put
\begin{equation}\label{eq:Pn}
 P_n(z)=\prod_{i=0}^{n-1}\prod_{j=0}^{n-1}
 G\left(z+\frac{i+j}{n}\right)
\end{equation}
and
\begin{equation}\label{eq:An}
 A_n(z)=n^{n^2z^2/2-nz}(2\pi)^{-(n^2-n)z/2}.
\end{equation}
A finite cancellation in the \(i\)-product gives
\begin{align}
 \frac{P_n(z+1/n)}{P_n(z)}
 &=\prod_{j=0}^{n-1}
   \frac{G(z+1+j/n)}{G(z+j/n)}\notag\\
 &=\prod_{j=0}^{n-1}\Gamma(z+j/n).
 \label{eq:Pshift}
\end{align}
Gauss' multiplication formula for Gamma is
\begin{equation}\label{eq:gauss-mult}
 \prod_{j=0}^{n-1}\Gamma(z+j/n)
 =(2\pi)^{(n-1)/2}n^{1/2-nz}\Gamma(nz).
\end{equation}
Also, direct subtraction of the exponents in \cref{eq:An} gives
\begin{equation}\label{eq:Ashift}
 \frac{A_n(z+1/n)}{A_n(z)}
 =n^{nz-1/2}(2\pi)^{-(n-1)/2}.
\end{equation}
Multiplying \cref{eq:Pshift,eq:gauss-mult,eq:Ashift},
\begin{equation}\label{eq:APshift}
 \frac{A_n(z+1/n)P_n(z+1/n)}{A_n(z)P_n(z)}=\Gamma(nz).
\end{equation}
The functional equation gives the same quotient for \(G(nz)\).  Hence
\begin{equation}\label{eq:periodic-ratio}
 Q_n(z):=\frac{G(nz)}{A_n(z)P_n(z)}
\end{equation}
is \(1/n\)-periodic wherever the quotient is defined; after cross
multiplication it extends through all removable points.

Apply the asymptotic expansion of \cref{sec:asymptotic} to every factor as
\(z\to+\infty\).  The coefficients of \(z^2\log z,z^2,z\log z,z,\log z\)
cancel.  The remaining constant is
\begin{equation}\label{eq:Kn}
 K(n)=\exp\bigl(-(n^2-1)\zeta'(-1)\bigr)
       n^{5/12}(2\pi)^{(n-1)/2}.
\end{equation}
Since
\begin{equation}\label{eq:A-zeta}
 \log A=\frac1{12}-\zeta'(-1),
\end{equation}
this is equivalently
\begin{equation}\label{eq:Kn-A}
 K(n)=\left(A\e^{-1/12}\right)^{n^2-1}
       n^{5/12}(2\pi)^{(n-1)/2}.
\end{equation}
Periodicity now makes the asymptotic determination exact: for fixed \(z\),
\(Q_n(z)=Q_n(z+m/n)\) for every integer \(m\), and the right side tends
effectively to \(K(n)\).

\begin{theorem}[Barnes multiplication]\label{thm:multiplication}
For every integer \(n\ge1\), with coherent branches for the powers,
\begin{equation}\label{eq:multiplication}
 G(nz)=K(n)n^{n^2z^2/2-nz}
 (2\pi)^{-(n^2-n)z/2}
 \prod_{i=0}^{n-1}\prod_{j=0}^{n-1}
 G\left(z+\frac{i+j}{n}\right).
\end{equation}
The branch-independent formal theorem is obtained by exponentiating the
quadratic logarithmic exponent along a declared path and cross multiplying all
Barnes factors.
\end{theorem}

\begin{proof}[Constructive modulus]
The finite shift calculation is exact.  Given a compact input box and precision
\(p\), choose an integer \(m\) large enough that every argument
\(z+m/n+(i+j)/n\) lies in the positive fixed sector used by the asymptotic
module and that the sum of the \(n^2+1\) explicit remainder bounds is at most
\(2^{-p-c}\).  Such an \(m\) is a primitive-recursive function of the box,
\(n\), the fixed truncation order and \(p\).  By exact periodicity, the desired
ratio at \(z\) equals the ratio at the shifted point.  The finite asymptotic
algebra reduces its ball to the constant \(K(n)\).  This proves every precision
instance.  The only use of a limit has thus been replaced by a checked finite
shift and a remainder inequality.
\end{proof}

\subsection{Duplication}

Set \(n=2\) in \cref{eq:multiplication}.  Since
\[
 P_2(x)=G(x)G(x+\tfrac12)^2G(x+1)
\]
and
\[
 K(2)=A^3\e^{-1/4}2^{5/12}(2\pi)^{1/2},
\]
solving for \(P_2(x)\) gives
\begin{equation}\label{eq:duplication}
 G(x)G(x+\tfrac12)^2G(x+1)
 =\e^{1/4}A^{-3}
 2^{-2x^2+3x-11/12}\pi^{x-1/2}G(2x).
\end{equation}
This is the Wikipedia duplication formula.  The powers of two are checked by
the exact exponent arithmetic
\[
 -2x^2+2x+x-\frac{5}{12}-\frac12
 =-2x^2+3x-\frac{11}{12}.
\]

\subsection{The value at one half}

Put \(x=1/2\) in \cref{eq:duplication}.  Since \(G(1)=1\) and
\(G(3/2)=\Gamma(1/2)G(1/2)=\sqrt\pi G(1/2)\),
\[
 \sqrt\pi\,G(1/2)^2=\e^{1/4}A^{-3}2^{1/12}.
\]
Positivity on the positive real axis selects the positive square root:
\begin{equation}\label{eq:Ghalf}
 G(1/2)=A^{-3/2}\pi^{-1/4}\e^{1/8}2^{1/24}.
\end{equation}
Two applications of the functional equation give
\begin{align}
 G(3/2)&=A^{-3/2}\pi^{1/4}\e^{1/8}2^{1/24},
 \label{eq:Gthreehalf}\\
 G(5/2)&=A^{-3/2}\pi^{3/4}\e^{1/8}2^{-23/24}.
 \label{eq:Gfivehalf}
\end{align}
These prove MathWorld equations (19)--(21).

All square roots here are positive real names.  Their iPRA proof uses lower and
upper dyadic bisection: from a positive input ball and a rational lower bound,
a bounded integer search returns the unique dyadic interval whose squared
endpoints enclose the input.

\subsection{Quarter values and Catalan's constant}

Define Catalan's constant by the alternating series
\begin{equation}\label{eq:Catalan}
 \mathcal K=\sum_{r=0}^{\infty}\frac{(-1)^r}{(2r+1)^2}.
\end{equation}
The alternating remainder is at most \((2N+3)^{-2}\), so \(\mathcal K\) has an
immediate primitive-recursive Cauchy modulus.  It is also
\(\Cl(\pi/2)\).

Let
\(a=\log G(1/4)\) and \(b=\log G(3/4)\).  Kinkelin reflection at \(z=1/4\)
gives, because \(G(5/4)=\Gamma(1/4)G(1/4)\),
\begin{equation}\label{eq:quarter-difference}
 b-a=\log\Gamma(1/4)-\frac18\log2-\frac14\log\pi
      +\frac{\mathcal K}{2\pi}.
\end{equation}
Here we used
\[
 \int_0^{1/4}\pi x\cot(\pi x)\dd x
 =\frac18\log2+\frac{\mathcal K}{2\pi}.
\]

Duplication at \(x=1/4\), followed by \cref{eq:Ghalf}, gives
\begin{equation}\label{eq:quarter-sum}
 a+b=\frac3{16}-\frac94\log A-\frac18\log2
 -\frac14\log\pi-\frac12\log\Gamma(1/4).
\end{equation}
Solving the two linear equations yields
\begin{equation}\label{eq:Gquarter}
 G(1/4)=A^{-9/8}\Gamma(1/4)^{-3/4}
 \exp\!\left(\frac3{32}-\frac{\mathcal K}{4\pi}\right),
\end{equation}
\begin{equation}\label{eq:Gthreequarter}
 G(3/4)=A^{-9/8}\Gamma(3/4)^{-1/4}
 \exp\!\left(\frac3{32}+\frac{\mathcal K}{4\pi}\right).
\end{equation}
Using
\(\Gamma(1/4)\Gamma(3/4)=\pi\sqrt2\), the second formula is equivalently
\begin{equation}\label{eq:Gthreequarter-alt}
 G(3/4)=A^{-9/8}2^{-1/8}\pi^{-1/4}
 \Gamma(1/4)^{1/4}
 \exp\!\left(\frac3{32}+\frac{\mathcal K}{4\pi}\right).
\end{equation}
Thus MathWorld's equation (18) is an alternative form of its preceding
\(G(3/4)\) value; it is not a third quarter argument.

The constructive proof does not solve a real linear system abstractly.  At
precision \(p\), it computes balls for the sum and difference with two guard
bits and applies the exact rational maps \((s,d)\mapsto((s-d)/2,(s+d)/2)\).
Exponentiation then uses a fixed positive-real evaluator.

\subsection{All odd half-integer values}

For \(k\ge2\), define
\begin{equation}\label{eq:ck}
 c_k=\prod_{i=1}^{k-2}
 \frac{2^i\Gamma(i+1/2)}{\sqrt\pi}
 =\prod_{i=1}^{k-2}(2i-1)!!.
\end{equation}
The second equality follows from
\(
\Gamma(i+1/2)=(2i)!\sqrt\pi/(4^ii!)
\), so \(c_k\) is an exact natural number.  Its primitive recursion is
\begin{equation}\label{eq:ck-rec}
 c_2=1,
 \qquad
 c_{k+1}=c_k(2k-3)!!.
\end{equation}

\begin{theorem}[Odd half-integer formula]\label{thm:odd-half}
For \(k>1\),
\begin{equation}\label{eq:odd-half}
 G(k-1/2)=
 c_k\,
 \frac{\pi^{(2k-3)/4}\e^{1/8}2^{1/24}}
      {2^{(k-1)(k-2)/2}A^{3/2}}.
\end{equation}
\end{theorem}

\begin{proof}
The case \(k=2\) is \cref{eq:Gthreehalf}.  Suppose the formula holds for
\(k\).  The functional equation gives
\[
 G(k+1/2)=\Gamma(k-1/2)G(k-1/2).
\]
Write
\[
 \Gamma(k-1/2)=\frac{(2k-3)!!}{2^{k-1}}\sqrt\pi.
\]
Multiplication by this factor changes \(c_k\) to \(c_{k+1}\), increases the
power of \(\pi\) by \(1/2\), and increases the denominator's power of two by
\(k-1\).  These are exactly the changes from \(k\) to \(k+1\) in
\cref{eq:odd-half}.  This is a quantifier-free induction once powers and
\(c_k\) are primitive-recursive symbols.
\end{proof}

\subsection{Formal dependency summary}

The multiplication theorem needs only the functional equation, Gauss
multiplication for Gamma, and the Barnes asymptotic constant.  Duplication is
its \(n=2\) instance.  The half values use duplication and recurrence.  The
quarter values use duplication, reflection and the single Clausen value
\(\Cl(\pi/2)=\mathcal K\).  Therefore the special-value module should be built
last; it contains almost no new analysis, only exact exponent arithmetic and
positive-real root/exponential operations.
